\documentclass[12pt,a4paper]{article}
\usepackage[utf8]{inputenc}
\usepackage[spanish]{babel}
\usepackage{hyperref}
\usepackage{geometry}
\geometry{margin=2.5cm}

\title{Práctica Next.js – Explorador de Pokémon}
\author{[Nome do estudante / Parella]}
\date{}

\begin{document}

\maketitle

\section*{Objetivo}

Desarrollar una aplicación web con \textbf{Next.js} que consuma la API pública \textbf{PokéAPI} y permita:

\begin{itemize}
    \item Listar todos los Pokémon en una página principal.
    \item Buscar Pokémon por nombre en tiempo real.
    \item Navegar a una página dinámica con el detalle de cada Pokémon.
    \item Volver a la página principal desde la vista de detalle.
\end{itemize}

La práctica tiene como objetivo trabajar:

\begin{itemize}
    \item Enrutado estático y dinámico en Next.js.
    \item Componentización.
    \item Fetching de datos desde una API REST.
    \item Manejo de estado.
    \item Props.
    \item Navegación con \texttt{Link} y/o \texttt{useRouter}.
    \item Estilado con CSS clásico.
\end{itemize}

\section*{API a utilizar}

\textbf{PokéAPI}

Endpoint base recomendado para listar Pokémon con su nombre y imagen:

\begin{verbatim}
https://pokeapi.co/api/v2/pokemon?limit=151
\end{verbatim}

Para obtener detalles de cada Pokémon:

\begin{verbatim}
https://pokeapi.co/api/v2/pokemon/{nombre}
\end{verbatim}

\section*{Requisitos funcionales}

\subsection*{1 Página Principal (/)}
Debe:

\begin{itemize}
    \item Mostrar todos los Pokémon obtenidos desde la API.
    \item Renderizar cada Pokémon mediante un componente reutilizable \texttt{<PokemonCard />}.
    \item Cada tarjeta debe mostrar:
    \begin{itemize}
        \item Imagen del Pokémon.
        \item Nombre.
    \end{itemize}
    \item Incluir un input de búsqueda que:
    \begin{itemize}
        \item Permita filtrar Pokémon por nombre.
        \item Filtre en tiempo real (sin recargar la página).
    \end{itemize}
\end{itemize}

\subsection*{2 Componente PokemonCard}
Debe:

\begin{itemize}
    \item Recibir por \texttt{props} los datos del Pokémon.
    \item Mostrar imagen y nombre.
    \item Ser clicable.
    \item Redirigir a la ruta dinámica correspondiente al Pokémon.
\end{itemize}

Ejemplo de ruta esperada:

\begin{verbatim}
/pokemon/pikachu
/pokemon/charmander
/pokemon/squirtle
\end{verbatim}

\subsection*{3 Página dinámica de Pokémon}

Ruta dinámica obligatoria:

\begin{verbatim}
/pokemon/[name]
\end{verbatim}

Debe:

\begin{itemize}
    \item Obtener la información del Pokémon seleccionado.
    \item Mostrar una página completa con al menos:
    \begin{itemize}
        \item Nombre
        \item Imagen
        \item Tipos
        \item Habilidades
        \item Peso y altura
    \end{itemize}
    \item Incluir un botón o enlace visible para volver a la página principal.
\end{itemize}

\section*{Estilado}

\begin{itemize}
    \item Se debe usar CSS clásico (CSS Modules o archivo global).
    \item Se valorará:
    \begin{itemize}
        \item Diseño limpio y legible.
        \item Buen uso de espaciado.
        \item Responsive básico.
        \item Efectos hover en tarjetas.
        \item Jerarquía visual clara.
    \end{itemize}
\end{itemize}

\section*{Requisitos técnicos}

\begin{itemize}
    \item Proyecto creado con \texttt{create-next-app}.
    \item Uso de rutas dinámicas.
    \item Componentes reutilizables.
    \item Uso correcto de \texttt{Link} o \texttt{useRouter}.
    \item Manejo correcto de estado para el buscador.
    \item Código limpio y organizado.
    \item Estructura de carpetas coherente.
\end{itemize}

\section*{Entrega}

\begin{itemize}
    \item Subir el proyecto a \textbf{GitHub}.
    \item Incluir un \texttt{README.md} con:
    \begin{itemize}
        \item Instrucciones para ejecutar el proyecto.
        \item Breve explicación de la estructura.
        \item Miembros participantes si se hace en pareja.
    \end{itemize}
\end{itemize}

\section*{Evaluación}

Se valorará especialmente:

\begin{itemize}
    \item Correcto funcionamiento de las rutas dinámicas.
    \item Calidad del código.
    \item Componentización.
    \item Buen uso de Next.js.
    \item Experiencia de usuario.
    \item Diseño y CSS.
    \item Limpieza del repositorio.
\end{itemize}

\end{document}